\section{Discussion}\label{sec:discussion}

\subsection{What occurrence-level validation adds}

Agreement varied far more across classes than across the tested span boundaries. Per-class F1 ranged from \num{0.481} to \num{0.974}, whereas micro F1 changed by \num{0.010} across the thresholds in Table~\ref{tab:iouSensitivity}. The pooled score therefore masks class-specific differences that matter when interpreting corpus flags. Occurrence-level matching exposes these differences because each match must align the class, the distinct event, and its source span.

The disagreement mechanisms in Table~\ref{tab:detectionErrors} help explain the class variation. Query-method semantics affected Implicit Columns and Poor Man's Search Engine. Inferred relationships and missing-value conventions affected Keyless Entry and Fear of the Unknown. Occurrence grouping and span boundaries caused further mismatches. The detector-informed revision raised micro F1 to \num{0.932}, showing how the estimate changes under an optimistic review of detector disagreements. The original-reference analysis remains primary.

The validation comparison supports a practical configuration choice. Zero-Shot prompting \nobreakcite{DBLP:journals/csur/LiuYFJHN23,DBLP:conf/nips/BrownMRSKDNSSAA20} with Claude Opus 4.5 \nobreakcite{opus45Card} combined the fastest Zero-Shot runtime with the lowest cost among the Claude configurations. Chain-of-Thought and Tree-of-Thought-inspired prompting \nobreakcite{DBLP:conf/nips/KojimaGRMI22,DBLP:conf/nips/Wei0SBIXCLZ22,tree-of-thought-prompting,DBLP:conf/nips/YaoYZS00N23} produced only a small observed F1 increase at higher cost and runtime. The single observed run per configuration supports the selection decision alone.

\subsection{What corpus normalisation adds}

The normalised results distinguish widespread detector output from repeated output within a file. The detector flags ID Required in almost half of generated-schema files and generally produces one flag in each flagged file. Implicit Columns flags occur in about one quarter of application/query files and produce nearly three flags per flagged file. Their raw totals therefore arise through different mechanisms, which a class ranking alone would hide.

Repository size also explains only part of the concentration. The highest-count decile contains about half of the relevant files and still has roughly one and a half times the density of the remaining repositories. Large repositories as a complete stratum have lower pooled density than small and medium repositories. These two results indicate that high raw counts reflect both source volume and repository-specific output composition rather than a general tendency for every large repository to have high density. Exact-content sensitivity preserves this interpretation, so byte-identical copies have little influence on the main pattern.

\subsection{Relation to existing detectors}

Existing detectors analyse different representations and output units. DbDeo and SQLInspect recover SQL text from host-language source before applying parsing and rules \nobreakcite{sharma2018smelly,nagy2017static,nagy2018sqlinspect}, and SQLCheck applies static and dynamic checks to executed queries \nobreakcite{DBLP:conf/sigmod/DintyalaNA20}. Prior LLM work evaluates smell recognition or classification in PL/SQL source \nobreakcite{de2025effectiveness}. This study instead analyses Java DSL expressions and generated schema classes and returns multiple class-labelled line spans per file.

A meaningful comparison should evaluate detectors on the same source files, class definitions, reference occurrences, and one-to-one matching rule. The priority comparison is a deterministic baseline for the seven-class jOOQ task. It could encode syntactically explicit patterns as API or AST rules \nobreakcite{DBLP:journals/tse/MohaGDM10} and reserve contextual cases for an LLM. Comparing these designs under one task would show how occurrence agreement, cost, and runtime change.

\subsection{Using source-fragment strata for corpus audit}

The recurring forms in Table~\ref{tab:apiManifestations} define practical strata for corpus validation. Implicit Columns flags concentrate in \texttt{selectFrom} and \texttt{select()} fragments, while Poor Man's Search Engine flags concentrate in \texttt{like(} fragments. An independent audit can sample flagged and unflagged uses within these strata, resolve call targets, and count relevant API uses to estimate class-specific errors and exposure-normalised associations.

The protocol could later assess transfer to other query builders and broader SQL-antipattern taxonomies \nobreakcite{alshemaimri2021survey,10.1007/978-3-031-09070-7_26,10.1007/978-3-031-35311-6_51,factor2014sql,karwin2026more} using representation-specific corpora and annotations.
